\clearpage
\SourceLeaf{228}\pdfbookmark[1]{Section Two.}{liber-secundus-sectio-secunda}
\begin{center}
{\large SECTION TWO.}\par\medskip
\emph{Determination of the orbit from four observations, of which only two are complete.}\par\bigskip
\hypertarget{native-228-article164}{}\pdfbookmark[2]{164.}{article-164}164.
\end{center}
At the very beginning of Book Two (Article 115), we stated that the problem treated in the preceding section applies only to orbits whose inclination neither vanishes nor is too small, and that the determination of slightly inclined orbits must necessarily be based on four observations. Four complete observations, however, are equivalent to eight equations, and there are only six unknowns, so the problem would be overdetermined. The latitudes (or declinations) must therefore be set aside in two of the observations so that the remaining data can be satisfied exactly. This gives rise to the problem to which this section is devoted. The solution we give here is applicable to slightly inclined orbits and can be used with equal success for orbits of any degree of inclination. Here, too, as in the problem of the preceding section, we must distinguish the case in which the approximate dimensions of the orbit are already known from the first determination of a hitherto entirely unknown orbit. We shall begin with the former.
